\section{Objective and Theorem Ladder}
\label{sec:v8-objective}
\label{sec:v8-theorem-ladder}

\subsection{The C-TreePO Objective}

Section~\ref{sec:v8-ctree-math} gives the structural goal: learn one
state map \(g\) whose node states can replace the manifesto spans they
represent for oracle-compatible readouts. We now write the loss that
pushes \(g\) toward that goal. The shape is two-channel. The root
channel asks the predicted document score to match the expert mean.
The local channel asks every realized state to behave like a valid
manifesto state in its own right.

Let \(L_{\mathrm{root}}(g)\) be the root-level loss, measured against
the available expert target. In the manifesto experiments,
\(L_{\mathrm{root}}\) measures the discrepancy between the predicted
root score and the expert-survey mean; the precise per-node agreement
metric used by the prompt-and-demonstration optimizer is given in
Section~\ref{sec:v8-benoit-tree-setup}.

For the local channel, attach one node loss to each \(v\in V(T)\).
The form of the loss depends on which law applies at \(v\). At a
leaf, the loss measures C1: does the leaf summary preserve the policy
evidence in the raw text? At a state passed back through \(g\), the
loss measures C2: does summarizing again give the same expert-facing
answer? At an internal merge, the loss measures C3: does merging the
two children preserve the joint policy evidence of the parent span?
If every check could be measured by the target oracle, the local-law
loss at node \(v\) would be
\[
  \ell_v^\ast
  =
  \ell_{\mathrm{law}}(v;f^\ast,g).
\]
The star marks the target quantity: the local-law loss as it would be
measured by the expert oracle, not by the cheaper proxy used during
dense training.

Let \(d(v)\) be the depth of node \(v\), with the root at depth \(0\),
and let \(\gamma\in[0,1]\) be a depth-discount weight. Each node enters
the local-law sum with weight \(\gamma^{d(v)}\), so the root counts
fully and deeper nodes count progressively less. \(\gamma=1\) treats
every node equally; smaller \(\gamma\) emphasizes nodes near the root
and damps deeper nodes. This is the same discount pattern used in
discounted-reward objectives, and the matching Lean formalization
carries \(\gamma\) through the local-law sum, so we report \(\gamma\)
as a declared knob even when \(\gamma=1\) is the operational default.
The oracle local-law channel is
\begin{equation}
\label{eq:v8-local-law-estimand}
  L_{\mathrm{law}}^\ast(g,T)
  =
  \sum_{v\in V(T)}\gamma^{d(v)}\ell_v^\ast .
\end{equation}
The theorem-facing objective is
\begin{equation}
\label{eq:v8-main-objective}
  J_\lambda^\ast(g)
  =
  (1-\lambda)L_{\mathrm{root}}(g)
  +
  \lambda L_{\mathrm{law}}^\ast(g,T),
  \qquad 0\le\lambda\le 1 .
\end{equation}
Here \(\lambda\) is the analyst-chosen local-law share. \(\lambda=0\)
trains on the document-level expert target alone, with no local-law
pressure on \(g\). \(\lambda=1\) trains on local laws alone and lets
the root drift. Intermediate values balance the two. The root term
names the manifesto score we ultimately report; the local term names
the state-validity pressure that makes the theorem relevant to the
fitted tree.

\subsection{Three Claim Stacks}

The objective sits inside three nested claim stacks. Each stack adds premises
and supports a different conclusion.

\begin{table}[t]
\centering
\small
\begin{tabular}{@{}p{0.24\linewidth}p{0.33\linewidth}p{0.34\linewidth}@{}}
\toprule
Stack & Added premise & Manifesto-facing conclusion \\
\midrule
Preservation & C1/C2/C3 hold on the realized C-Tree for \(f^\ast\) &
Node and root states can replace their raw manifesto spans for
oracle-compatible readouts. \\
Optimization & Preferences or losses depend on the input only through the
preserved oracle interface &
Training on valid root summaries has the same population target as training
on full manifestos. \\
Certificate & Approximate local-law distortion is bounded or estimated from
sampled nodes with logged propensities &
Approximate preservation becomes a finite-sample bound on downstream
objective drift. \\
\bottomrule
\end{tabular}
\caption{\textbf{The theorem ladder.} Each tier keeps the Manifesto target
explicit: the oracle is the policy measurement the analyst wants to preserve.}
\label{tab:v8-theorem-ladder}
\end{table}

\begin{table}[t]
\centering
\small
\begin{tabular}{@{}lcccccccc@{}}
\toprule
Result & C1 & C2 & C3 & OC & SP & CL & Lip & FT \\
\midrule
Thm.~\ref{thm:v8-root-preservation} & \checkmark & \checkmark & \checkmark & & & & & \checkmark \\
Cor.~\ref{cor:v8-preferences-pass-to-root} & \checkmark & \checkmark & \checkmark & \checkmark & & & & \checkmark \\
Prop.~\ref{prop:v8-manifesto-preference-equivalence} & \checkmark & \checkmark & \checkmark & \checkmark & & & & \checkmark \\
Prop.~\ref{prop:v10-population-gap} & \checkmark & \checkmark & \checkmark & \checkmark & & & \checkmark & \checkmark \\
Prop.~\ref{prop:v10-corrected-loss-unbiased} & & & & & \checkmark & & & \checkmark \\
Prop.~\ref{prop:v10-ht-unbiased} & & & & & \checkmark & & & \checkmark \\
Prop.~\ref{prop:v10-finite-sample-certificate} & & & & \checkmark & \checkmark & \checkmark & \checkmark & \checkmark \\
\bottomrule
\end{tabular}
\caption{\textbf{Assumption-to-claim crosswalk.} Each row marks the
assumptions that the corresponding formal result invokes. C1, C2, C3 are
the local laws of Definition~\ref{def:v8-local-laws}; OC is
oracle-compatibility of the downstream readout or loss; SP is strict
positivity \(\pi_i>0\) of logged inclusion probabilities; CL is a
bounded clipping radius on inverse-propensity weights; Lip is the
method transport constant \(C_{\mathrm{meth}}\) (a Lipschitz or
equivalent regularity condition on the downstream method); FT is
Assumption~\ref{ass:v8-fixed-tree} (fixed realized tree).}
\label{tab:v10-assumption-claim}
\end{table}

The three stacks build on each other. Preservation is the cleanest
claim and assumes exact local laws. Optimization keeps the exact-laws
premise and adds a compatibility condition on the downstream training
object. Certificate relaxes exact laws to a sampled bound estimated
from the realized tree. We take them in order.

The preservation stack is Theorem~\ref{thm:v8-root-preservation} and
Corollary~\ref{cor:v8-preferences-pass-to-root}. It is the
local-to-global claim: if every realized state remains in the same
\(f^\ast\)-fiber as the span it represents, then the root state
remains in the same fiber as the whole manifesto.

The optimization stack adds an oracle-compatibility premise for the
downstream training object. The manifesto experiments below use a
regression target, where this premise holds automatically:
\(L_{\mathrm{root}}\) sees the manifesto only through its predicted
score. The premise becomes load-bearing when the framework is applied
to preference-learning or RL-style downstream tasks. For pairwise
preference learning, a comparison rule is compatible when replacing a
manifesto by an equivalent root state does not change the preferred
answer. This covers DPO-style pairwise objectives
\citep{BradleyTerry1952,RafailovEtAl2023DPO} and GRPO-style group
objectives only when the prompt, response generator, policy, ranker,
or reward channel depends on the manifesto through the preserved
policy oracle. A preference for formatting, rhetoric, or style
outside the policy target needs a wider oracle interface.

\begin{proposition}[Manifesto Preference Objective Equivalence]
\label{prop:v8-manifesto-preference-equivalence}
Fix a population of manifestos and a realized C-Tree construction for each
manifesto. If C1/C2/C3 hold exactly for \(f^\ast\), and if the downstream
preference loss is oracle-compatible in the input argument, then replacing each
manifesto by its C-Tree root state leaves the population preference objective
unchanged.
\end{proposition}

\begin{proof}[Proof sketch]
Theorem~\ref{thm:v8-root-preservation} gives
\(z_{\mathrm{root}}\sim_\ast x\) for every manifesto in the
population. Oracle-compatibility of the preference loss lifts
\(\sim_\ast\) to identity of the loss value at every realization.
Taking expectations over the manifesto population preserves
equality. Full proof in Appendix~\ref{app:v8-full-proofs};
machine-checked as
\texttt{paper\_preference\_stack\_same\_argmin}, with method-specific
exact-metric surfaces \texttt{dpo\_exact\_metric},
\texttt{grpo\_pl\_exact\_metric}, \texttt{grpo\_rl\_exact\_metric} in
\texttt{FormalProofs/OPT/MainTheorems.lean} and
\texttt{FormalProofs/OPT/PreferenceLearning.lean}.
\end{proof}

The proposition says that root summaries can be used as training inputs only
for preference tasks whose input dependence is already captured by the
manifesto policy target. Exact local laws handle the compression step; the
oracle-compatibility premise handles the downstream objective.

\subsection{Approximate Preservation}

Exact local laws are the clean theorem endpoint. The deployed LLM
setting is approximate. The realized state at each node differs from
the raw span by some amount, and the realized score at the root
differs from the expert mean by some amount. Three quantities organize
that approximation.

\(G_{\mathrm{meth}}\) is the drift in whatever downstream quantity the
analyst is reporting (a population correlation, a population MSE, or a
preference-task population gap). This is the error the analyst cares
about.

\(\Delta_R\) is the average distortion between a produced state and
the raw span it represents, measured by the oracle. The subscript
\(R\) marks the represented-span comparison.

\(C_{\mathrm{meth}}\) is a method-specific transport constant. It
captures how much the downstream quantity can move when the
representation moves: a small \(C_{\mathrm{meth}}\) means a smooth
downstream method that tolerates small representation distortion well,
and a large \(C_{\mathrm{meth}}\) means a sharper method where small
representation errors translate into bigger reported drift.

\begin{proposition}[Population Gap Bound]\label{prop:v10-population-gap}
Fix a downstream method with transport constant \(C_{\mathrm{meth}}\)
and represented-span distortion \(\Delta_R\). The drift in the
reported quantity is bounded by
\begin{equation}
\label{eq:v8-population-gap}
  |G_{\mathrm{meth}}|
  \le
  C_{\mathrm{meth}}\,\Delta_R .
\end{equation}
\end{proposition}

\begin{proof}[Proof sketch]
\(C_{\mathrm{meth}}\) is exactly the Lipschitz constant linking
representation distortion in the oracle metric to the downstream
quantity. Apply the Lipschitz inequality and take expectations under
the manifesto population. Full proof in
Appendix~\ref{app:v8-full-proofs}; the regression case uses readout
Lipschitzness; DPO uses
\texttt{dpo\_loss\_pointwise\_lipschitz}; the unified preference
bound is \texttt{unified\_preference\_gap\_bounded} in
\texttt{FormalProofs/OPT/PreferenceBounds.lean} and
\texttt{FormalProofs/OPT/MainTheorems.lean}.
\end{proof}

For the manifesto regression case, \(C_{\mathrm{meth}}\) is the
Lipschitz constant of the score readout: how far the predicted score
can move when the input state moves a little. For DPO,
\(C_{\mathrm{meth}}\) packages a policy-Lipschitz condition on the
log-probability ratios. For group-ranking or RL-style objectives,
\(C_{\mathrm{meth}}\) packages the ranker, reward, clipping, and KL
terms. The constant belongs to the downstream training method; the
C-Tree audit supplies \(\Delta_R\).

Section~\ref{sec:v8-audit-certificate} develops the operational route
for \(\Delta_R\): sample realized manifesto tree nodes and use a
Horvitz--Thompson estimator with logged inclusion
probabilities.\footnote{A separate representation route bounds
\(\Delta_R\) from function-class approximation properties of \(g\); we
do not develop that route in the main text because the operational
route is what carries the deployed manifesto measurement. The
function-class formulation is in
Appendix~\ref{app:v8-neural-operator-route}.}
Appendix~\ref{app:v8-full-proofs} gives the proof details for the
preservation, preference-transport, and certificate statements used
in this ladder.
